\chapter*{واژه‌نامه}
\addcontentsline{toc}{chapter}{واژه‌نامه}
\markboth{واژه‌نامه}{واژه‌نامه}

این واژه‌نامه برای یکسان‌ماندن معنای اصطلاحات در سراسر کتاب تنظیم شده است. تعریف‌ها کوتاه‌اند تا هنگام مطالعه بتوان به آن‌ها رجوع کرد؛ توضیح مفصل هر مفهوم در فصل مربوط به خودش آمده است.

\section*{داده، اطلاعات و بازنمایی}

\begin{description}
  \item[داده (\lr{Data})] تفاوت یا محتوایی که در یک زمینه‌ی مشخص ثبت، دریافت یا بررسی می‌شود. داده به‌تنهایی لزوماً به معنای فهم یا دانش نیست.
  \item[اطلاعات (\lr{Information})] محتوایی که در یک چارچوب تفسیر می‌شود و می‌تواند عدم‌قطعیت یا ناآگاهی را کاهش دهد. در نظریه‌ی شانون، مقدار اطلاعات معنای فنی و احتمالاتی مشخصی دارد.
  \item[دانش (\lr{Knowledge})] ساختاری از اطلاعاتِ فهم‌شده و مرتبط که می‌تواند در توضیح، استدلال یا عمل به کار رود.
  \item[خرد (\lr{Wisdom})] توانایی داوری درباره‌ی کاربرد دانش، با توجه به زمینه، پیامدها و ارزش‌ها.
  \item[بازنمایی (\lr{Representation})] رابطه یا سازوکاری که به‌واسطه‌ی آن چیزی می‌تواند در یک سامانه به‌جای چیز دیگری قرار گیرد یا آن را نشان دهد.
  \item[نشانه (\lr{Sign})] چیزی که در یک رابطه‌ی تفسیری به چیزی دیگر اشاره می‌کند یا آن را بازنمایی می‌کند.
  \item[نماد (\lr{Symbol})] نشانه‌ای که رابطه‌اش با موضوع عمدتاً بر قرارداد، قاعده یا عادت تفسیری تکیه دارد.
  \item[حامل فیزیکی (\lr{Physical Carrier})] وضعیت یا فرایند مادی‌ای که یک الگو یا تفاوتِ قابل تشخیص در آن تحقق می‌یابد و می‌تواند در یک سامانه نقش اطلاعاتی داشته باشد.
\end{description}

\section*{ارتباط و نظریه‌ی اطلاعات}

\begin{description}
  \item[پیام (\lr{Message})] محتوایی که قرار است از منبع به مقصد منتقل شود؛ پیام را نباید با سیگنال فیزیکیِ حامل آن یکی دانست.
  \item[منبع (\lr{Source})] سامانه‌ای که پیام را تولید می‌کند یا یکی از پیام‌های ممکن را انتخاب می‌کند.
  \item[رمزگذاری (\lr{Encoding})] تبدیل پیام به رشته‌ای از نمادها یا شکل دیگری که برای انتقال یا ذخیره‌سازی مناسب باشد.
  \item[رمزگشایی (\lr{Decoding})] بازسازی پیام یا نمادهای آن از روی چیزی که دریافت شده است.
  \item[کانال (\lr{Channel})] مسیری که سیگنال از طریق آن از فرستنده به گیرنده می‌رسد؛ کانال می‌تواند مادی، نوری، صوتی، زیستی یا ذخیره‌ساز باشد.
  \item[نویز (\lr{Noise})] تغییرات یا اختلال‌هایی که تشخیص سیگنال یا بازسازی پیام را دشوار می‌کنند.
  \item[احتمال (\lr{Probability})] عددی میان صفر و یک که در یک مدل، میزان امکان رخ‌دادن یک نتیجه را توصیف می‌کند.
  \item[عدم‌قطعیت (\lr{Uncertainty})] ناآگاهی درباره‌ی اینکه کدام نتیجه از میان چند امکان رخ خواهد داد؛ در نظریه‌ی اطلاعات می‌توان آن را با توزیع احتمال کمّی کرد.
  \item[خوداطلاعات (\lr{Self-information})] مقدار غافلگیریِ یک نتیجه‌ی مشخص، که در نظریه‌ی شانون با \(-\log_2 p\) بیان می‌شود.
  \item[آنتروپی شانون (\lr{Shannon Entropy})] میانگین خوداطلاعات نتیجه‌های یک متغیر تصادفی؛ معیاری از عدم‌قطعیت توزیع احتمال، نه سنجه‌ای برای معنای پیام.
  \item[ظرفیت کانال (\lr{Channel Capacity})] بیشینه‌ی نرخ اطلاعاتی که با توجه به مدل کانال می‌توان به‌طور قابل اعتماد منتقل کرد.
  \item[بیت (\lr{Bit})] واحدی برای مقدار اطلاعات که با انتخاب یکی از دو حالت هم‌احتمال متناظر است؛ در محاسبات اطلاعاتی، مقادیر بزرگ‌تر بر پایه‌ی لگاریتم دودویی بیان می‌شوند.
  \item[پیچیدگی کولموگروف (\lr{Kolmogorov Complexity})] طول کوتاه‌ترین برنامه‌ای که یک رشته را روی یک ماشین مرجع تولید می‌کند؛ این مفهوم با آنتروپی شانون یکسان نیست.
\end{description}

\section*{نشانه‌ها، طبیعت و فیزیک}

\begin{description}
  \item[شمایل (\lr{Icon})] نشانه‌ای که به‌واسطه‌ی شباهت یا ویژگی مشترک با موضوع خود ارتباط پیدا می‌کند.
  \item[نمایه (\lr{Index})] نشانه‌ای که با موضوع خود پیوند واقعی، علّی یا مکانی دارد؛ مانند دودی که می‌تواند نشانه‌ی آتش باشد.
  \item[نشانه‌ی طبیعی (\lr{Natural Sign})] اثر یا وضعیتی طبیعی که بدون قصد ارتباطی می‌تواند شواهدی درباره‌ی رویداد یا وضعیت دیگری فراهم کند.
  \item[جسم نامتناوب (\lr{Aperiodic Solid})] تعبیر شرودینگر برای ساختاری فرضی از ماده‌ی وراثتی که برخلاف بلورهای ساده، آرایشی پیچیده و تکرارنشونده دارد.
  \item[فیزیکالیسم (\lr{Physicalism})] دیدگاهی فلسفی که واقعیت را در نهایت فیزیکی می‌داند یا معتقد است هر واقعیت به واقعیت فیزیکی وابسته است.
  \item[حالت منطقی (\lr{Logical State})] مقدار انتزاعی‌ای مانند صفر یا یک که در یک سامانه‌ی منطقی تفسیر می‌شود.
  \item[حالت فیزیکی (\lr{Physical State})] وضعیت واقعی یک سامانه‌ی مادی که می‌تواند با کمیت‌هایی مانند ولتاژ، بار، میدان یا آرایش مولکولی توصیف شود.
  \item[شیطان ماکسول (\lr{Maxwell's Demon})] آزمایش فکری‌ای که رابطه‌ی میان اطلاعات، اندازه‌گیری، انتخاب و قانون دوم ترمودینامیک را به پرسش می‌کشد.
  \item[اصل لاندائر (\lr{Landauer's Principle})] اصلی که برای پاک‌کردن منطقیِ برگشت‌ناپذیر اطلاعات، حد پایین هزینه‌ی گرمایی را تعیین می‌کند.
  \item[\lr{DNA}] مولکولی که توالی واحدهای شیمیایی آن در ذخیره و انتقال اطلاعات وراثتی نقش دارد.
\end{description}

\section*{سیگنال و الکتریسیته}

\begin{description}
  \item[سیگنال (\lr{Signal})] تغییر یا رویداد فیزیکیِ قابل تشخیصی که در یک سامانه‌ی مشخص می‌تواند حامل پیام یا اطلاعات باشد.
  \item[سیگنال آنالوگ (\lr{Analog Signal})] سیگنالی که یک کمیت آن می‌تواند در یک بازه به‌صورت پیوسته تغییر کند.
  \item[سیگنال دیجیتال (\lr{Digital Signal})] سیگنالی که تغییرات فیزیکی آن به مجموعه‌ای از حالت‌ها یا نمادهای گسسته تعبیر می‌شود.
  \item[نمونه‌برداری (\lr{Sampling})] اندازه‌گیری سیگنال در لحظه‌های زمانی گسسته.
  \item[علیاسینگ (\lr{Aliasing})] پدیده‌ای که در آن نمونه‌برداری ناکافی باعث می‌شود مؤلفه‌ها یا سیگنال‌های متفاوت در داده‌ی نمونه‌برداری‌شده به‌اشتباه قابل تشخیص نباشند.
  \item[کوانتیزه‌سازی (\lr{Quantization})] نگاشت مقدارهای پیوسته یا دقیق به مجموعه‌ای محدود از سطح‌های گسسته.
  \item[بار الکتریکی (\lr{Electric Charge})] ویژگی فیزیکی ماده که در برهم‌کنش الکترومغناطیسی نقش دارد و با کولن اندازه‌گیری می‌شود.
  \item[پتانسیل الکتریکی (\lr{Electric Potential})] انرژی پتانسیل الکتریکی به‌ازای واحد بار در یک نقطه، نسبت به مرجع انتخاب‌شده.
  \item[اختلاف پتانسیل (\lr{Potential Difference})] تفاوت پتانسیل الکتریکی میان دو نقطه؛ کمیتی که با ولت اندازه‌گیری می‌شود.
  \item[ولتاژ (\lr{Voltage})] نام رایج اختلاف پتانسیل الکتریکی میان دو نقطه.
  \item[آستانه‌ی منطقی (\lr{Logic Threshold})] مرزی در یک مدار یا سامانه‌ی دیجیتال که برای تشخیص ناحیه‌های متناظر با حالت‌های منطقی به کار می‌رود.
\end{description}
